\chapter{Methodology}
\label{experimental_design}
This chapter outlines the design decisions underlying the empirical evaluation of this study. It describes the selected datasets and models, the evaluation metrics, the validation strategy and the hyperparameter tuning approach that together form the methodological foundation for the experiments.

% =============================================================
\section{Datasets}

\section{Models}
\label{design_models}


Building on the tree-based approach, \gls{xgb} was chosen as a second mid-complexity model, with the key distinction that it also supports \gls{gpu}-accelerated training. This makes it possible to directly compare the energy consumption and training time of \gls{cpu}- and \gls{gpu}-based tree ensemble methods under identical conditions, which is central to \rqref{3}. The suitability of \gls{gpu}-accelerated \gls{xgb} for datasets at this scale has been demonstrated by Mitchell and Frank, who show that the histogram-based \gls{gpu} implementation can process a dataset of 10 million instances with 28 features entirely within \gls{gpu} memory \cite{mitchellAcceleratingXGBoostAlgorithm2017}.

The \gls{mlp} was selected as the most complex model, representing the \gls{nn} category. Its highly configurable architecture, particularly the number of hidden layers and neurons per layer, makes it well suited for the architectural complexity experiments conducted in this study, where these parameters are systematically varied to assess their impact on predictive performance and energy consumption. Like \gls{lr}, the \gls{mlp} also requires standard scaling, which is applied identically.




\gls{rf} is the only model excluded from one dataset combination: due to prohibitive training runtimes at 11 million instances, it is not evaluated on HIGGS in the main benchmark. It is, however, included in the dataset scaling experiment for subset sizes up to 500{,}000 instances, where runtimes remain manageable. This provides data on the point at which \gls{rf} becomes computationally impractical relative to the other models. All other model--dataset combinations are included in the main benchmark.

\section{Evaluation Metrics}
\label{design_metrics}

To answer the research questions of this study, two categories of metrics were employed: classification performance and ecological cost.
\Gls{wf1} was selected to measure predictive performance, as it provides a standardized basis for comparing models across datasets of varying class distributions.
Energy consumption in \gls{kwh} and \gls{co2eq} emissions in \gls{gco2eq} were chosen as the primary resource metrics, directly reflecting the ecological cost of model training.

Training time in seconds was included as a secondary resource metric, as it captures computational demand independently of hardware-specific power draw.
Furthermore, inference time was measured across all experimental setups to account for the operational footprint. In large-scale deployment scenarios, the cumulative energy demand of repeated predictions can often rival or even exceed the initial training costs \cite{desislavovTrendsAIInference2023}.
Together, these metrics enable a direct assessment of whether gains in predictive performance justify the associated increase in training cost.

Inspired by the efficiency framing proposed in \citeauthor{schwartzGreenAI2020} \cite{schwartzGreenAI2020}, a custom composite metric was developed for this study, named \gls{cf1}. \Cref{eq:cf1} defines this metric as the ratio of training \gls{co2eq} emissions (in \gls{mgco2eq}) to the \gls{wf1} (in percent):
\begin{equation}
\text{\gls{cf1}} = \frac{C}{\gls{wf1}}
\label{eq:cf1}
\end{equation}
where $C$ denotes \gls{co2eq} emissions in \gls{mgco2eq}. Milligrams are chosen as the unit because training emissions in this study range from roughly 0.01\,\gls{gco2eq} to 90\,\gls{gco2eq}, which divided by \gls{wf1} scores near 70--99\,\% yields \gls{cf1} values in the range of approximately 0.1 to 1{,}000, which are human-readable integers that would become unwieldy fractions in grams. Lower values indicate a more ecologically efficient model: less carbon expenditure per unit of predictive performance. Unlike normalization-based composite metrics, \gls{cf1} is an absolute measure: its value is independent of which other models appear in the same comparison, making it meaningful even when a single model is evaluated in isolation.

It should be noted that $C$ is inherently location- and time-dependent, as \gls{co2eq} emissions are tied to the carbon intensity of the local electricity grid, which varies both by region and over time~\cite{wrightEfficiencyNotEnough2025}. Consequently, \gls{cf1} scores reflect the ecological efficiency of the specific experimental setup used in this study and may not generalize directly to other energy infrastructures or time periods.

Emissions are tracked separately for three distinct phases: hyperparameter tuning, model training and inference. In all three cases, data loading and preprocessing are deliberately excluded from the measurement interval, as their cost scales with dataset size rather than model architecture and would therefore introduce noise into cross-model comparisons. This ensures that reported emissions reflect only the algorithmic cost of each phase.

\section{Validation Strategy}
\label{design_validation}

To obtain an unbiased estimate of generalization performance, the evaluation methodology employs a strict separation between model development and final testing. A shuffled, stratified 80/20 train--test split is applied to each dataset prior to any modeling steps. The resulting training partition is used exclusively for \gls{hpo} and model fitting, while the held-out test set is reserved solely for the final evaluation.

Within the \gls{hpo} phase, \gls{cv} is used as the internal evaluation strategy on the training partition. Specifically, shuffled Stratified K-Fold \gls{cv} with $k=5$ folds is applied, a value chosen for its established balance between computational cost and estimation stability. Random shuffling is applied before fold construction to ensure that the order of instances in the original dataset does not influence fold composition. Stratified splitting then ensures that each fold preserves the original class distribution, which is particularly relevant because two of the three datasets exhibit class imbalance. This approach is directly coherent with the choice of \gls{wf1} as the primary metric and evaluation objective.

Final model evaluation is performed by training a single model on the full training partition and scoring it on the unseen test set. By strictly withholding the test data from the \gls{cv} loop, this protocol ensures that the hyperparameters are not implicitly optimized for the test set, thereby preventing over-fitting during the model selection phase.

To ensure reproducibility, a fixed random state was used for all splits, \gls{cv} folds, model initialization, and data shuffling. The specific value and its technical implementation are described in \Cref{implementation_sw_env}.

A practical limitation applies to the Wine dataset: an 80/20 split of its 178 instances leaves approximately 36 test samples, which yields a noisier point estimate than the thousands or millions of test instances available for Credit and HIGGS. This is a structural constraint of the dataset size rather than a methodological choice, and the split ratio is kept uniform across all datasets to preserve comparability.

\section{Applied HPO}
\label{design_hyperparameter}
Hyperparameter tuning was applied to all models except \gls{lr}. Without tuning, default parameters tend to favor certain model families over others, potentially skewing both performance and training cost comparisons \cite{bagnallUseDefaultParameter2017}. \gls{lr} is deliberately excluded: its role as a lightweight linear baseline depends on its simplicity and tuning it would undermine that function.

\gls{bo}-based methods require substantially fewer evaluations than grid or random search to reach competitive configurations, which directly reduces the energy cost of the tuning process itself. Because of that, \texttt{Optuna} was selected as the optimization framework, using \gls{tpe} as its search strategy. Each run was configured to execute 40 trials per model--dataset combination, a budget chosen to balance search thoroughness against ecological overhead. All trials operate exclusively on the training partition and tuning was performed separately for each dataset. The optimization objective is \gls{wf1} in all cases.

A limitation of the fixed 40-trial budget is that it does not account for differences in search space dimensionality across models. \gls{rf} exposes five hyperparameters and \gls{xgb} six, both in well-defined, low-dimensional spaces that \gls{tpe} can navigate reliably within 40 evaluations. The \gls{mlp} search space is larger and structurally more complex: it includes a discrete depth parameter (\texttt{n\_layers}, ranging from 1 to 4) that conditionally activates up to four additional layer-width dimensions, plus continuous parameters for dropout rate, learning rate and batch size, yielding an effective dimensionality of five to eight depending on the sampled depth. Conditional search spaces of this kind require more evaluations for \gls{tpe} to build an accurate surrogate model. In addition, each \gls{mlp} trial on HIGGS involves training a \gls{nn} and is substantially more expensive than a comparable tree-based trial, so the 40-trial budget represents a greater practical constraint for the \gls{mlp} than for \gls{rf} or \gls{xgb}. As a result, the \gls{mlp} may not have reached the same degree of optimization as the tree-based models and the reported \gls{mlp} performance figures should be interpreted as a lower bound on what a more exhaustive search might achieve. Increasing the trial budget selectively for the \gls{mlp} would, however, break the comparability constraint: allocating more search effort to one model than to others would confound any observed performance difference with tuning thoroughness rather than isolating it as a property of the algorithm itself. The uniform budget is therefore a deliberate design choice to keep the search effort symmetric across all models.


